\documentclass[11pt,a4paper]{article}

% --- Paket ---
\usepackage[utf8]{inputenc}
\usepackage[T1]{fontenc}
\usepackage[swedish]{babel}
\usepackage{amsmath, amssymb, amsthm}
\usepackage{mathtools}
\usepackage{geometry}
\geometry{margin=2.5cm}
\usepackage{parskip}
\usepackage{xcolor}
\usepackage{hyperref}
\hypersetup{
    colorlinks=true,
    linkcolor=blue!50!black,
    urlcolor=blue!50!black
}

% --- Egna kommandon ---
\newcommand{\R}{\mathbb{R}}
\newcommand{\E}{\mathbb{E}}
\newcommand{\Var}{\operatorname{Var}}
\newcommand{\dd}{\mathrm{d}}
\DeclareMathOperator{\Cov}{Cov}

% --- Satsmiljöer ---
\theoremstyle{plain}
\newtheorem{theorem}{Sats}
\newtheorem{lemma}[theorem]{Lemma}

\theoremstyle{definition}
\newtheorem{definition}[theorem]{Definition}

\theoremstyle{remark}
\newtheorem*{remark}{Anmärkning}

% --- Titel ---
\title{Härledning av den analytiska lösningen \\ till Ornstein--Uhlenbeck-processen}
\author{Projektarbete SF1693}
\date{\today}

\begin{document}

\maketitle

\section{Problemformulering}

Vi betraktar Ornstein--Uhlenbeck-processen (OU-processen), definierad av den stokastiska differentialekvationen
\begin{equation}
    \dd X_t = \theta(\mu - X_t)\,\dd t + \sigma\,\dd W_t, \qquad X_0 = x_0,
    \label{eq:ou-sde}
\end{equation}
där $\theta > 0$ är återgångshastigheten, $\mu \in \R$ är långtidsmedelvärdet, $\sigma > 0$ är volatiliteten, och $\{W_t\}_{t \geq 0}$ är en standard Wienerprocess. Målet är att härleda en explicit lösning $X_t$, samt bestämma dess fördelning.

\section{Förberedande räkneregler}

Innan vi löser \eqref{eq:ou-sde} sammanfattar vi de verktyg från Itôs kalkyl som kommer att användas.

\begin{lemma}[Itôs lemma]
\label{lem:ito}
Antag att $X_t$ uppfyller en SDE av formen $\dd X_t = b_t\,\dd t + s_t\,\dd W_t$, och låt $f \in C^{2,1}(\R \times [0,\infty))$. Då gäller
\begin{equation}
    \dd f(X_t, t) = \frac{\partial f}{\partial t}\,\dd t + \frac{\partial f}{\partial x}\,\dd X_t + \frac{1}{2}\frac{\partial^2 f}{\partial x^2}(\dd X_t)^2,
    \label{eq:ito}
\end{equation}
där $(\dd X_t)^2$ beräknas med räknereglerna
\[
    \dd t \cdot \dd t = 0, \qquad \dd t \cdot \dd W_t = 0, \qquad \dd W_t \cdot \dd W_t = \dd t.
\]
\end{lemma}

\begin{lemma}[Itô-isometrin]
\label{lem:isometri}
För en deterministisk funktion $f \in L^2([0,t])$ gäller
\[
    \E\!\left[\left(\int_0^t f(s)\,\dd W_s\right)^{\!2}\,\right] = \int_0^t f(s)^2\,\dd s.
\]
\end{lemma}

\begin{lemma}[Gaussianitet av Wiener-integraler]
\label{lem:gauss}
Om $f \colon [0,t] \to \R$ är deterministisk och kvadratiskt integrerbar, så är
\[
    \int_0^t f(s)\,\dd W_s \;\sim\; \mathcal{N}\!\left(0,\;\int_0^t f(s)^2\,\dd s\right).
\]
\end{lemma}

\begin{proof}[Bevisskiss till lemma~\ref{lem:gauss}]
Itô-integralen definieras som ett $L^2$-gränsvärde av Riemann-summor
\[
    S_n = \sum_{i=0}^{n-1} f(s_i)\bigl(W_{s_{i+1}} - W_{s_i}\bigr).
\]
Eftersom Wiener-inkrementen är oberoende och normalfördelade, och $f(s_i)$ är deterministisk, är varje $S_n$ en linjärkombination av oberoende gaussiska variabler --- alltså själv gaussisk. Klassen av gaussiska fördelningar är sluten under $L^2$-konvergens, så även gränsvärdet är gaussiskt. Medelvärde och varians följer av linjäritet respektive Itô-isometrin.
\end{proof}

\section{Härledning av lösningen}

\subsection*{Steg 1: Integrerande faktor}

Vi vill eliminera den linjära driften i \eqref{eq:ou-sde} och prövar samma grepp som för en linjär ODE: integrerande faktor. Betrakta först det deterministiska fallet $\dot X + \theta X = \theta\mu$. Multiplikation med $e^{\theta t}$ samlar vänsterledet till en ren tidsderivata,
\[
    \frac{\dd}{\dd t}\!\left[e^{\theta t}(X - \mu)\right]
    = e^{\theta t}\bigl(\dot X + \theta(X - \mu)\bigr) = 0,
\]
så storheten $e^{\theta t}(X - \mu)$ bevaras av ODE-flödet. Detta motiverar att i det stokastiska fallet studera den motsvarande processen
\begin{equation}
    Y_t \coloneqq e^{\theta t}(X_t - \mu).
    \label{eq:Y-def}
\end{equation}
Frågan är om Itôs lemma ger en lika ren reduktion av $\dd Y_t$ --- det vill säga, om driften försvinner även stokastiskt. Vi kontrollerar det i nästa steg.

\subsection*{Steg 2: Tillämpa Itôs lemma}

Låt $f(x,t) = e^{\theta t}(x - \mu)$, så att $Y_t = f(X_t, t)$. De partiella derivatorna är
\[
    \frac{\partial f}{\partial t} = \theta e^{\theta t}(x - \mu), \qquad
    \frac{\partial f}{\partial x} = e^{\theta t}, \qquad
    \frac{\partial^2 f}{\partial x^2} = 0.
\]
Att andraderivatan är noll är avgörande: eftersom $f$ är linjär i $x$ försvinner Itô-korrektionen. Lemma~\ref{lem:ito} ger
\begin{equation}
    \dd Y_t = \theta e^{\theta t}(X_t - \mu)\,\dd t + e^{\theta t}\,\dd X_t.
    \label{eq:dY-step1}
\end{equation}
Sätt in SDE:n \eqref{eq:ou-sde} för $\dd X_t$:
\begin{align}
    \dd Y_t &= \theta e^{\theta t}(X_t - \mu)\,\dd t + e^{\theta t}\bigl[\theta(\mu - X_t)\,\dd t + \sigma\,\dd W_t\bigr] \nonumber \\
    &= \underbrace{\theta e^{\theta t}(X_t - \mu)\,\dd t - \theta e^{\theta t}(X_t - \mu)\,\dd t}_{= \,0} + \sigma e^{\theta t}\,\dd W_t.
    \label{eq:dY-cancel}
\end{align}
Drifttermerna kancellerar exakt --- precis som väntat. Vi har nu
\begin{equation}
    \boxed{\dd Y_t = \sigma e^{\theta t}\,\dd W_t.}
    \label{eq:dY-final}
\end{equation}

\subsection*{Steg 3: Integrera}

Ekvation \eqref{eq:dY-final} är drivlös med deterministisk diffusionskoefficient. Integrera direkt från $0$ till $t$:
\begin{equation}
    Y_t - Y_0 = \int_0^t \sigma e^{\theta s}\,\dd W_s = \sigma \int_0^t e^{\theta s}\,\dd W_s.
    \label{eq:Y-integral}
\end{equation}
Eftersom $Y_0 = e^{0}(X_0 - \mu) = x_0 - \mu$ följer
\begin{equation}
    Y_t = (x_0 - \mu) + \sigma \int_0^t e^{\theta s}\,\dd W_s.
    \label{eq:Y-explicit}
\end{equation}

\subsection*{Steg 4: Lös ut $X_t$}

Från definitionen \eqref{eq:Y-def} har vi $X_t = \mu + e^{-\theta t} Y_t$. Insättning av \eqref{eq:Y-explicit} ger
\begin{align}
    X_t &= \mu + e^{-\theta t}\!\left[(x_0 - \mu) + \sigma \int_0^t e^{\theta s}\,\dd W_s\right] \nonumber \\
    &= \mu + (x_0 - \mu)e^{-\theta t} + \sigma e^{-\theta t}\int_0^t e^{\theta s}\,\dd W_s.
\end{align}
Eftersom $e^{-\theta t}$ är oberoende av integrationsvariabeln $s$ kan den föras in i integralen:
\begin{equation}
    \boxed{X_t = \mu + (x_0 - \mu)e^{-\theta t} + \sigma \int_0^t e^{-\theta(t-s)}\,\dd W_s.}
    \label{eq:ou-solution}
\end{equation}
Detta är OU-processens analytiska lösning.

\section{Fördelningen av $X_t$}

\subsection*{Medelvärde}

Itô-integralen av en deterministisk integrand är en martingal som startar i noll, så
\[
    \E\!\left[\int_0^t e^{-\theta(t-s)}\,\dd W_s\right] = 0.
\]
Av \eqref{eq:ou-solution} följer
\begin{equation}
    \E[X_t] = \mu + (x_0 - \mu)e^{-\theta t}.
    \label{eq:mean}
\end{equation}
När $t \to \infty$ konvergerar medelvärdet exponentiellt mot $\mu$ med hastighet $\theta$.

\subsection*{Varians}

Endast Itô-integralen i \eqref{eq:ou-solution} bidrar till variansen, eftersom övriga termer är deterministiska. Eftersom integralen har medelvärde noll gäller
\begin{align}
    \Var(X_t) &= \sigma^2 \,\E\!\left[\left(\int_0^t e^{-\theta(t-s)}\,\dd W_s\right)^{\!2}\,\right] \nonumber \\
    &= \sigma^2 \int_0^t e^{-2\theta(t-s)}\,\dd s \quad\text{(Itô-isometrin)}.
\end{align}
Substitutionen $u = t - s$, $\dd u = -\dd s$, ger
\[
    \int_0^t e^{-2\theta(t-s)}\,\dd s = \int_0^t e^{-2\theta u}\,\dd u = \left[-\frac{e^{-2\theta u}}{2\theta}\right]_0^t = \frac{1 - e^{-2\theta t}}{2\theta}.
\]
Således
\begin{equation}
    \boxed{\Var(X_t) = \frac{\sigma^2}{2\theta}\bigl(1 - e^{-2\theta t}\bigr).}
    \label{eq:variance}
\end{equation}

\subsection*{Fördelning}

Itô-integralen i \eqref{eq:ou-solution} är gaussisk enligt lemma~\ref{lem:gauss}, eftersom integranden $e^{-\theta(t-s)}$ är deterministisk. Resten av uttrycket för $X_t$ är deterministiskt, så $X_t$ är en affin transformation av en gaussisk variabel --- och därmed själv gaussisk. Med medelvärdet \eqref{eq:mean} och variansen \eqref{eq:variance} har vi visat följande resultat.

\begin{theorem}[Fördelningen för OU-processen]
\label{thm:ou}
Lösningen $X_t$ till \eqref{eq:ou-sde} med startvärde $X_0 = x_0$ är normalfördelad enligt
\begin{equation}
    X_t \mid X_0 = x_0 \;\sim\; \mathcal{N}\!\left(\mu + (x_0 - \mu)e^{-\theta t},\; \frac{\sigma^2}{2\theta}\bigl(1 - e^{-2\theta t}\bigr)\right).
\end{equation}
\end{theorem}

\section{Stationär fördelning}

När $t \to \infty$ ger \eqref{eq:mean} och \eqref{eq:variance}
\[
    \E[X_t] \to \mu, \qquad \Var(X_t) \to \frac{\sigma^2}{2\theta}.
\]
Den stationära fördelningen är därmed
\begin{equation}
    \pi = \mathcal{N}\!\left(\mu,\; \frac{\sigma^2}{2\theta}\right),
    \label{eq:stationary}
\end{equation}
oberoende av startvärdet $x_0$. Detta uttrycker OU-processens \emph{ergodicitet}: oavsett var processen startar konvergerar dess fördelning mot $\pi$ med hastighet bestämd av $\theta$.

\section{Sammanfattning}

Genom integrerande faktorn $e^{\theta t}$ och transformationen $Y_t = e^{\theta t}(X_t - \mu)$, som eliminerar driften via Itôs lemma, har vi härlett att OU-processen ges av
\[
    X_t = \mu + (x_0 - \mu)e^{-\theta t} + \sigma \int_0^t e^{-\theta(t-s)}\,\dd W_s,
\]
och att $X_t \sim \mathcal{N}\bigl(m(t), v(t)\bigr)$ med
\[
    m(t) = \mu + (x_0 - \mu)e^{-\theta t}, \qquad v(t) = \frac{\sigma^2}{2\theta}\bigl(1 - e^{-2\theta t}\bigr).
\]
Den stationära fördelningen är $\mathcal{N}\bigl(\mu, \sigma^2/(2\theta)\bigr)$. Dessa formler utgör referenslösningen som den numeriska lösaren av motsvarande Fokker--Planck-ekvation kommer att valideras mot.

\end{document}